% p.65 — الباب التاسع والعشرون في الاصطلاحات والتركيبات (تتمة: التركيبات)
\mspage{65}
% Layout: four columns = two pairs (label | text), right pair first, continuing the
% تركيب examples of p.64: the first row completes example (7).  Circled example numbers
% written as (n).  Printed page number ۶۵ at top left.  Layout: the two pairs are set side
% by side (right pair first) in one `tables` row, as in the lithograph.
\begin{tables}
\begin{sarf}{|r|r|}
جملة صفة & ضرب غلامه عمرا جملة صفة\\
 & (8) جَاءَ رَجُلٌ الَّذِيْ غُلَامُهُ ضَرَبَ عَمْرًا\\
موصول & الَّذِيْ موصول\\
جملة صله & غلامه ضرب عمرا جملة صلة\\
 & (9) اِنَّ زَيْدًا قَائِمٌ\\
حرف مشبه بالفعل & انّ حرف مشبه بالفعل\\
اسم & زيدا اسم ان\\
خبر & قائم خبر ان\\
 & (10) كَانَ زَيْدٌ قَائِمًا\\
فعل ناقص & كان فعل ناقص\\
اسم & زيد اسم كان\\
خبر & قائمًا خبر كان\\
 & (11) كَادَ زَيْدٌ أَنْ يَخْرُجَ\\
فعل مقاربة & كاد فعل مقاربة\\
فاعل مقاربة & زيد فاعل مقاربة\\
خبر مقاربة & ان يخرج خبر مقاربة\\
 & (12) نِعْمَ رَجُلٌ زَيْدٌ\\
فعل مدح & نعم فعل مدح\\
فاعل مدح & رَجُلٌ فاعل مدح\\
مخصوص بالمدح & زيد مخصوص بالمدح\\
 & (13) بِئْسَ رَجُلٌ زَيْدٌ\\
\end{sarf}
\begin{sarf}{|r|r|}
فعل ذم & بئس فعل ذم\\
فاعل ذم & رجل فاعل ذم\\
مخصوص بالذم & زيد مخصوص بالذم\\
 & (14) ضَرَبَ زَيْدٌ عَمْرًا الْيَوْمَ ضَرْبًا\\
فعل & ضرب فعل\\
فاعل & زيد فاعل\\
مفعول به & عمرا مفعول به\\
مفعول فيه & اليوم مفعول فيه\\
مفعول مطلق & ضربا مفعول مطلق\\
 & (15) ضَرَبَ زَيْدٌ عَمْرًا تَأْدِيْبًا\\
مفعول له & تأديبا مفعول له\\
 & (16) اِسْتَوَى الْمَاءُ وَالْخَشَبَةَ\\
مفعول معه & الخشبة مفعول معه\\
معيه & واو معية\\
 & (17) عَلِمَ زَيْدٌ بَكْرًا صَالِحًا\\
فعل يقين & عَلِمَ فعل يقين\\
مفعول اول & بكرا مفعول اول\\
مفعول ثانى & صالحا مفعول ثانى\\
 & (18) ظَنَّ زَيْدٌ بَكْرًا صَالِحًا\\
فعل شك & ظن فعل شك\\
 & (19) سَرَقَ زَيْدٌ الْقَلَمَ\\
\end{sarf}
\end{tables}
